\chapter{Performance Optimized Digital Logic}\label{chapt:perf_digital_logic}
\localtoc

\section{Comparison of \texttt{when} and \texttt{MuxCase} in Chisel} \label{sec:muxcase}

The choice between \texttt{when} and \texttt{MuxCase} in Chisel depends on the specific design scenario and synthesis goals. 

\subsection{Behavioral and Structural Implications}
\texttt{when} is more natural and intuitive for conditions that resemble software-like \texttt{if-else} logic. It typically generates cascaded conditional logic (e.g., \texttt{if-else} structures) in the hardware, where each condition is checked sequentially. On the other hand, \texttt{MuxCase} represents a priority-free selection logic, a multiplexer, where all conditions are evaluated simultaneously. This structure often synthesizes directly into a multiplexer tree, which can result in better performance and reduced latency for certain designs.

\subsection{Synthesis Implications}
\textbf{Timing and Performance:}  
For timing, \texttt{when} introduces propagation delays due to the cascading nature of sequential conditions. In contrast, \texttt{MuxCase} generates parallel logic, potentially yielding better timing performance as conditions are evaluated independently.

\textbf{Area:}  
The \texttt{when} construct can produce more compact hardware for sequential conditions since unused conditions are pruned early. However, \texttt{MuxCase} may generate larger hardware if many mutually exclusive conditions are evaluated together because all conditions are checked in parallel.

\textbf{Readability:}  
\texttt{when} is preferable for sequential, human-readable conditions with natural priority, while \texttt{MuxCase} is better suited for a concise representation of mutually exclusive cases.

\subsection{Specific Considerations}
\textbf{Mutually Exclusive Conditions:}  
If conditions are naturally exclusive, \texttt{MuxCase} is often a better fit. For example:
\[
\texttt{val result = MuxCase(0.U, Seq(}
\]
\[
\quad (\texttt{opcode === ADD} \rightarrow \texttt{(a + b),}
\]
\[
\quad (\texttt{opcode === SUB} \rightarrow \texttt{(a - b))}
\]
\texttt{))}
This translates to a 3-to-1 multiplexer in hardware.

\textbf{Priority Logic:}  
When conditions have a natural priority (e.g., an \texttt{else} case that catches all unhandled conditions), \texttt{when} is preferable. For example:
\[
\texttt{when(opcode === ADD) \{}
\]
\[
\quad \texttt{result := a + b}
\]
\[
\texttt{\}.elsewhen(opcode === SUB) \{}
\]
\[
\quad \texttt{result := a - b}
\]
\[
\texttt{\}.otherwise \{}
\]
\[
\quad \texttt{result := 0.U}
\]
\texttt{\}}  
This logic represents a chain of conditional checks in hardware.

\subsection{Switch Synthesis in Chisel}
In Chisel, a \texttt{switch} statement is a construct used for multi-way branching based on a condition or value. The hardware synthesized from a \texttt{switch} statement depends on the conditions and the synthesis toolchain. Generally, it can result in one of the following hardware structures:

A \textbf{multiplexer tree} is generated if the \texttt{switch} conditions are mutually exclusive, with each case corresponding to a distinct value of the selector. For instance:

\begin{lstlisting}[language=Scala]
val result = Wire(UInt(8.W)) 
switch(selector) { 
 is("b00".U) { result := a } 
 is("b01".U) { result := b } 
 is("b10".U) { result := c } 
 is("b11".U) { result := d } 
 }
\end{lstlisting}

This should synthesize into a 4-to-1 multiplexer controlled by the value of \texttt{selector}.

When the \texttt{switch} includes overlapping or non-mutually-exclusive conditions, \textbf{priority logic} is typically generated. Priority logic evaluates cases sequentially, giving precedence to earlier conditions. For wide selectors or more complex logic, synthesis tools may generate \textbf{decoder logic} to translate the selector into one-hot or prioritized enable signals that control the corresponding cases.

\subsection{Comparison with \texttt{MuxCase} and \texttt{when}}
\texttt{MuxCase} is more explicit and directly creates a multiplexer, making it compact and efficient for mutually exclusive conditions. \texttt{when} constructs explicitly generate priority logic. In contrast, \texttt{switch} can represent either multiplexers or priority logic depending on the use case, providing a flexible but less predictable synthesis.

\subsubsection{Limitations of \texttt{switch}}
While \texttt{switch} provides a clean and readable way to handle multi-way branching, it is not always recommended for large or complex selectors. For wide selectors, the generated decoder logic can result in significant area and timing overhead. Additionally, overlapping or non-mutually-exclusive cases can inadvertently lead to priority chains that increase latency and complicate timing analysis. 


%**********************************
\section{Carry Lookahead Adder}
%*********************************
Carry Look-Ahead Adders (CLAs) are designed to overcome the performance bottlenecks of ripple-carry adders by reducing the dependency on sequential carry propagation. Instead of waiting for each bit's carry to be computed, the CLA computes carries in parallel using the generate (\texttt{G}) and propagate (\texttt{P}) signals for each bit. These signals allow the adder to determine whether a carry will be generated or propagated at each stage. Using these signals, the carry for each bit position can be calculated directly through a set of Boolean equations, significantly reducing the delay to logarithmic time relative to the adder's width. This design is particularly effective for high-speed arithmetic operations in processors and digital signal processing, where minimizing latency is important. However, the parallel computation and additional logic circuitry result in increased hardware complexity compared to simpler adders.

%******************************
\section{Carry Select Adder}
%******************************
Carry Select Adders (CSAs) improve upon the delay of ripple-carry adders by precomputing the sum and carry outputs for both possible carry-in values (0 and 1) in parallel. Each stage of the adder computes two possible results based on these assumptions, and a multiplexer is used to select the correct result once the carry-in is determined. This approach significantly reduces the propagation delay associated with sequential carry computation while maintaining a relatively straightforward design. Although CSAs require additional hardware resources for the parallel computations and multiplexers, they achieve a good trade-off between speed and area, making them well-suited for medium-sized arithmetic units where both latency and complexity are important considerations.

%*****************************
\section{Carry Save Adder}
%*****************************
Carry-save adders are specialized adders used primarily in multiplication. Unlike conventional adders, which propagate carries through the stages immediately, CSAs represent the result as two separate numbers: a partial sum and a carry vector. These two outputs are then combined in a subsequent addition step, often using a more efficient adder such as a carry-lookahead adder. By deferring the carry propagation, CSAs significantly reduce the delay for intermediate steps, making them highly efficient for scenarios where multiple additions must be performed simultaneously. Examples include the Wallace tree \cite{wallace1964} and Dadda tree \cite{dadda1965}. This efficiency is achieved at the cost of additional logic to manage the separate sum and carry outputs.






%*****************************
\section{Kogge-Stone Adder} \label{sec:kogge-stone-adder}
%*****************************
Parallel prefix adders are a class of high-performance adders designed to optimize the computation of carry signals in binary addition, significantly reducing delay compared to simpler adder designs like ripple-carry adders. They achieve this by dividing the addition process into three distinct stages: generating initial propagate and generate signals for each bit, computing intermediate carries through a tree-like prefix structure, and finally calculating the sum based on these carry signals. The tree structure allows the carry signals to be computed in parallel, with a logarithmic time complexity relative to the number of bits, making these adders highly efficient for large bit-width operations. These adders are widely used in modern high-speed computing systems and processors.

The Kogge-Stone Adder is a parallel prefix adder known for its high performance and suitability for large bit-width arithmetic operations. It minimizes the propagation delay using a tree-like structure to compute carry signals in a logarithmic number of stages relative to the bit width. The design uses generate and propagate signals to recursively compute carries in parallel, ensuring that the carry for any given bit is available after just $\log_2(n)$ stages, where $n$ is the bit width. While this structure significantly enhances speed, it requires substantial wiring and additional resources, making it resource-intensive compared to simpler designs like ripple-carry adders. Due to its efficiency in reducing delay, the Kogge-Stone Adder is frequently employed in high-performance processors and digital systems requiring fast arithmetic computations  \cite{kogge1973parallel}.


%%%%%%%%%%%%%%%%%%%%%%%%%%%%%%%%%%%%%%%%%%%%%%%%%%%%%%%%%%%%%%%%%%%%%%%%%%%%%
\section{Multi-Port Register Files}
%%%%%%%%%%%%%%%%%%%%%%%%%%%%%%%%%%%%%%%%%%%%%%%%%%%%%%%%%%%%%%%%%%%%%%%%%%%%%
A multi-port register file is an important component in high-performance processors, particularly in superscalar and out-of-order architectures, where multiple instructions execute concurrently. The design of such a register file requires consideration of factors including port scalability, read and write access, timing constraints, area efficiency, and power consumption. 

%****************************************************
\subsection{Port Requirements and Scalability}
%****************************************************
The number of read and write ports directly impacts the complexity and feasibility of the register file. In a superscalar machine with \( W \)-wide issue capability, a register file must typically support at least \( 2W \) read ports (for source operands) and \( W \) write ports (for destination results). For example, a 4-issue processor may require 8 read ports and 4 write ports per cycle.

As the number of ports increases, the cost of multiplexing, routing, and driving signals grows non-linearly. Fully combinational read access requires a large multiplexer per port, which increases propagation delay and power dissipation. This leads to a fundamental trade-off: more ports improve throughput but worsen area, power, and timing constraints.

%****************************************************
\subsection{Fully Combinational vs. Pipelined Access}
%****************************************************
A register file may use either fully combinational reads (single-cycle access) or pipelined reads (multi-cycle access). In fully combinational designs, every read operation produces a result in the same cycle, but this requires large fan-in multiplexers, leading to long critical paths. This approach is typically used for small, low-latency register files such as general-purpose registers in simple CPU cores.

In contrast, pipelined register files introduce a one-cycle read delay, allowing for better timing closure at high clock frequencies. This technique is commonly used in deeply pipelined superscalar processors to balance speed and complexity.


%****************************************************
\subsection{Multi-Porting Techniques}
%****************************************************
Since SRAMs do not natively support multiple read and write ports, multi-porting is achieved through three primary techniques: multi-banking, register replication, and bypass networks.

\subsubsection{Banked Register Files}
A \textbf{banked register file} partitions the register file into multiple independent SRAM banks, allowing concurrent accesses as long as no two accesses target the same bank. The bank assignment can be \textit{Static}, where specific registers are assigned to predetermined banks, or \textit{Dynamic}, where banking conflicts are resolved through remapping.

Banking is highly efficient but introduces potential bank conflicts, requiring additional logic to stall or redistribute requests.

\subsubsection{Register File Replication}
\textbf{Register replication} duplicates the entire register file across multiple SRAM blocks, with each block providing one read port per cycle. A 4-issue processor requiring 8 read ports might use 4 copies of the register file, each servicing 2 reads. While this approach eliminates read conflicts, it significantly increases area and power consumption.

\subsubsection{Bypass Networks and Shadow Registers}
In deeply pipelined designs, many register reads can be bypassed instead of accessing the main register file. A bypass network forwards the result of recently written values directly to dependent instructions, reducing pressure on the register file. Additionally, a shadow register file (implemented using Register) can store recently written values before committing them to an SRAM-based backing store.

%****************************************************
\subsection{Write-Back Arbitration}
%****************************************************
When multiple instructions write back in the same cycle, a **write port contention problem** occurs. To resolve this \textbf{Write port prioritization} assigns priority levels to different sources, such as integer ALUs, floating-point units, or memory loads. \textbf{Write combining} merges multiple writes into a single cycle if they target the same register. \textbf{Staggered write-back} spreads writes across consecutive cycles to avoid conflicts.

The write arbitration logic ensures that high-priority instructions complete their write-back first, preventing stalls in the execution pipeline.

%****************************************************
\subsection{Power and Area Considerations}
%****************************************************
Register files consume significant power due to frequent access by multiple functional units. To optimize power. \textbf{Clock gating} disables unused portions of the register file to save dynamic power.  \textbf{Read clustering} groups read ports together based on access patterns, reducing unnecessary accesses.
\textbf{Voltage scaling} lowers power consumption in lower-speed execution units.

Area optimizations include SRAM-based storage, narrow-width optimizations (storing frequently used 8-bit and 16-bit values in smaller memory cells), and hierarchical register files that separate frequently accessed registers from less frequently used ones.

%****************************************************
\subsection{Integration with Pipeline and Execution Units}
%****************************************************
A register file must interface with execution units while maintaining pipeline efficiency. The integration involves \textit{Read scheduling} to ensure that operands are available before execution. \textit{Write-back coordination} to update the register file without conflicts. \textit{Forwarding mechanisms} to prevent stalls by using bypass paths for recently computed results.

The placement of the register file in the processor pipeline also influences performance. Some designs use a centralized register file, while others distribute smaller register files across execution clusters to reduce wire delay.

%****************************************************
\subsection{Comparison of Multi-Porting Strategies}
%****************************************************
A comparison of multi-porting strategies is shown in the following table:

{\footnotesize
\begin{center}
\begin{tabular}{|l|c|c|c|}
\hline
\textbf{Method} & \textbf{Latency} & \textbf{Scalability} & \textbf{Area Efficiency} \\
\hline
\texttt{Vec[Reg]} (FF-based) & Zero-cycle reads & Poor for large sizes & High area and power \\
Multi-Banked SRAM & 1-cycle read delay & Good, requires conflict resolution & Efficient for large files \\
Register Replication & 1-cycle read delay & Scalable but increases power & High area cost \\
Bypass Networks & Zero-cycle forwarding & Enhances performance & Increases control  complexity \\
\hline
\end{tabular}
\end{center}
}

\ifthenelse{\boolean{chiselversion}}{
%************************************************
\subsection{Parameterized 2R/1W SRAM-based Register File} \label{subsec:regfile_2R1W_SRAM}
%**************************************************
\includechiselcode[caption={Parameterized 2R/1W SRAM-based Register File (Chisel: RegFile2R1WSRAM.scala)}, label={lst:chisel_2R1W_SRAM_regfile}, firstline=22, firstnumber=22, lastline=999]{\chiselSourceDir/memory/RegFile2R1WSRAM.scala}

Listing \ref{lst:chisel_2R1W_SRAM_regfile} shows \texttt{RegFile2R1WSRAM} that implements a register file with two read ports and one write port. It is built using a synchronous memory (\texttt{SyncReadMem}), ensuring that all memory operations are synchronized with the clock signal. The register file stores \texttt{depth} registers, where each register is \texttt{width} bits wide. By default, both \texttt{depth} and \texttt{width} are set to 32, meaning the register file contains 32 registers, each capable of holding a 32-bit value.

To address individual registers within the register file, the code computes the number of address bits required using the function \texttt{log2Ceil(depth)}. This function determines the minimum number of bits necessary to represent an index into a structure of size \texttt{depth}. The computed address width is then used to define input and output signals in the module's interface.

The module interface consists of several input and output ports. The inputs include \texttt{src1} and \texttt{src2}, which specify the addresses of the registers to be read, and \texttt{dst1}, which specifies the address of the register to be written. The \texttt{wen} signal acts as a write enable, determining whether a write operation should take place. The \texttt{dst1data} input carries the data to be written to the register file. The outputs \texttt{src1data} and \texttt{src2data} provide the data read from the registers at addresses \texttt{src1} and \texttt{src2}, respectively.

The register file is implemented using a synchronous memory block \texttt{mem}, instantiated with \texttt{SyncReadMem}. This memory is designed for synchronous reads, meaning that data is retrieved one clock cycle after the corresponding address is provided. The module reads data from the memory using \texttt{mem.read(io.src1)} and \texttt{mem.read(io.src2)}. Since the read operation is synchronous, the outputs must be registered to ensure they align correctly with the processor pipeline. This is achieved using \texttt{RegNext}, which registers the read outputs so that they are available one cycle later. If no previous value exists, the outputs default to zero.

The write operation is controlled by the \texttt{wen} signal. When \texttt{wen} is asserted, the module writes the value of \texttt{dst1data} into the memory at the address specified by \texttt{dst1}. This ensures that the write operation is executed only when explicitly enabled. Because the processor pipeline ensures that write operations do not coincide with reads in the same cycle, there is no need for additional bypass logic to resolve read-after-write hazards.



%************************************************
\subsection{Parameterized Multithreaded 2R/1W SRAM-based Register File} \label{subsec:regfile_MT_2R1W_SRAM}
%**************************************************


\includechiselcode[caption={Parameterized Multithreaded 2R/1W SRAM-based Register File (Chisel: RegFileMT2R1WSRAM.scala)}, label={lst:chisel_MT_2R1W_SRAM_regfile}, firstline=29, firstnumber=29, lastline=999]{\chiselSourceDir/memory/RegFileMT2R1WSRAM.scala}

Listing \ref{lst:chisel_MT_2R1W_SRAM_regfile} \texttt{RegFileMT2R1WSRAM} implements a barrel-threaded multithreaded register file with two read ports and one write port. This design utilizes a synchronous memory (\texttt{SyncReadMem}) to store register values while supporting multiple independent threads. Each thread maintains a separate register space, and a thread identifier (\texttt{threadID}) is used to distinguish between the register sets. Instead of storing registers separately for each thread, the module implements a unified SRAM structure, where the total memory depth is equal to the product of the number of threads and the number of registers per thread.

To compute memory addresses for accessing registers, the module concatenates the \texttt{threadID} with the register index. This concatenation effectively partitions the memory into per-thread sections, ensuring that each thread has exclusive access to its own registers. The number of bits required to represent a thread identifier is determined using \texttt{log2Ceil(threads)}, while the number of bits needed to index registers within a thread is computed using \texttt{log2Ceil(depth)}. The resulting memory depth, given by \texttt{threads * depth}, defines the total storage required for the register file.

The module's input and output interface consists of several ports that facilitate register file operations. The \texttt{threadID} input determines the active thread, and all read and write accesses are performed within the address space of this thread. The read ports \texttt{src1} and \texttt{src2} specify the register indices to be read, while their corresponding outputs \texttt{src1data} and \texttt{src2data} return the stored values after a one-cycle delay. The write port \texttt{dst1} specifies the register index for writing data, and \texttt{dst1data} contains the value to be written. The write operation is controlled by the \texttt{wen} signal, which enables writes when asserted.

To perform memory operations, the module first computes the effective memory addresses by concatenating the \texttt{threadID} with the respective register indices. These computed addresses are used for both read and write operations. The read operation is synchronous, meaning that data is not available immediately after providing the address but appears in the next clock cycle. To align with the processor pipeline, the module registers the read outputs using \texttt{RegNext}, ensuring that the read values are properly synchronized with the computation stages.

For write operations, when \texttt{wen} is asserted, the module writes the value of \texttt{dst1data} into memory at the computed effective address \texttt{Cat(io.threadID, io.dst1)}. This guarantees that each thread writes only to its own register space without interfering with other threads. The use of a unified memory structure simplifies the implementation while allowing independent register access for each thread.



%************************************************
%\subsection{Parameterized Multiport Register File}
%**************************************************
%\includechiselcode[caption={Parameterized Multi R/W port Register File (Chisel: RegFile.scala)}, label={lst:chisel_multiport_regfile}, firstline=9, firstnumber=9, lastline=999]{\chiselSourceDir/memory/RegFile.scala}

%Listing \ref{lst:chisel_multiport_regfile} 

}% End Chisel


%****************************************************
\subsection{Conclusion}
%****************************************************
The design of a multi-port register file is a complex trade-off between latency, power, area, and access bandwidth. Small, low-latency register files are best implemented using \texttt{Vec[Reg]}, while large-scale register files require SRAM-based architectures with multi-porting techniques such as banking, replication, and bypassing. A well-optimized register file ensures efficient execution of instructions while maintaining balance between scalability and implementation cost.


%*************************************
\section{Mealy vs Moore Automata}
%*************************************
Mealy and Moore automata are two fundamental models for designing finite state machines, and they differ primarily in the way outputs are generated. In a Mealy automaton, the output is a function of both the current state and the current input. This means that the output can change immediately in response to an input change, potentially resulting in faster response times. However, this tight coupling may also lead to glitches if the input changes asynchronously relative to the state transitions.

In contrast, a Moore automaton generates outputs solely based on the current state, independent of the immediate input. As a result, the output in a Moore machine only changes when the state changes. This can simplify the design and make the system more robust, since outputs are synchronized with state transitions, but it may also introduce a delay in output response compared to a Mealy design.

Another key difference is the state complexity. Mealy machines often require fewer states to implement the same functionality because they leverage the input signal directly to determine outputs. Moore machines, while potentially more straightforward in terms of output logic, might require additional states to achieve equivalent behavior since each state must uniquely determine the output.

In summary, the choice between a Mealy and a Moore automaton involves trade-offs between response speed, design simplicity, and state complexity. A Mealy design is advantageous when quick reaction to inputs is essential, whereas a Moore design may be preferred for its stability and easier analysis, especially in environments where synchronized outputs are critical.